% CCSS 7.SP.B.3 — Comparing Two Populations Visually
\section{Comparing Two Populations Visually}

\begin{learningGoals}
\begin{itemize}
    \item Compare two data distributions by visual overlap
    \item Measure difference between centers and variability
\end{itemize}
\end{learningGoals}

\begin{conceptBox}{Visual Comparison}
Compare distributions by looking at their graphs and overlap.
\end{conceptBox}

\begin{workedExample}{Comparing Heights}
Class A: mean $150$ cm, Class B: mean $155$ cm. Both have similar spread. Which is taller?
\answerHighlight{Class B is taller on average, but overlap means some students in A are taller than some in B.}
\end{workedExample}

\tipBox{Look at both the center (mean/median) and spread (range, MAD) when comparing.}

\begin{practiceBox}{Comparing Populations}
\resetProblems

\prob Two classes have means $80$ and $85$. Which is higher?
\answerExplain{Second class}{The mean of $85$ is higher than $80$.}

\prob If two groups have overlapping ranges, what does this mean?
\answerExplain{Some members of each group are similar}{Overlap means not all members are different.}

\prob Why is it important to look at variability?
\answerExplain{Shows how spread out the data is}{Variability helps understand differences beyond just the mean.}

\prob If Group A has mean $70$, Group B has mean $75$, but Group A has a larger range, what does this mean?
\answerExplain{Group A is more spread out}{A larger range means more variability.}

\prob What is mean absolute deviation (MAD)?
\answerExplain{Average distance from the mean}{MAD measures variability in a data set.}

\end{practiceBox}
